% Generated draft: requires mathematical and source review.
\documentclass[11pt]{article}
\usepackage[margin=25mm]{geometry}
\usepackage{fontspec}
\setmainfont{Latin Modern Roman}
\usepackage{amsmath,amssymb,mathtools,unicode-math}
\setmathfont{Latin Modern Math}
\usepackage{enumitem,needspace,xurl,hyperref}
\setlength{\parindent}{0pt}
\setlength{\parskip}{5pt}
\setlength{\emergencystretch}{3em}
\newcommand{\N}{\mathbb{N}}
\newcommand{\Z}{\mathbb{Z}}
\newcommand{\Q}{\mathbb{Q}}
\newcommand{\R}{\mathbb{R}}
\newcommand{\C}{\mathbb{C}}
\newcommand{\F}{\mathbb{F}}
\newcommand{\A}{\mathbb{A}}
\newcommand{\PP}{\mathbb{P}}
\newcommand{\E}{\mathbb{E}}
\newcommand{\eps}{\varepsilon}
\newcommand{\dd}{\,\mathrm{d}}
\begin{document}
\section*{UnsolvedMath ID: 800817}
\textbf{Source problem code:} catalogue entry
\subsection{Definitions and mathematical statement}
Let \(\mathscr{D}_\mu^s(M)\) denote the source's group of volume-preserving diffeomorphisms of Sobolev regularity \(H^s\) on \(M\), with identity \(e\). The geometric setting, the volume form \(\mu\), the range of \(s\), and any boundary conditions must be specified before this notation determines the intended geodesic problem. In the usual geometric hydrodynamics interpretation, a Riemannian metric \(g\) on \(M\) induces the right-invariant \(L^2\) metric whose value at the identity is
\[
\langle u,w\rangle_{L^2}=\int_M g(u,w)\,\mathrm d\mu,
\qquad u,w\in T_e\mathscr{D}_\mu^s(M).
\]
This is an interpretation of the source notation; the supplied record does not state the hypotheses ensuring a smooth geodesic flow and exponential map.

Where the geodesic exponential map \(\operatorname{Exp}_e\) exists and is differentiable, write
\[
\gamma_v(t)=\operatorname{Exp}_e(tv),\qquad \gamma_v(0)=e.
\]
For a positive time \(t\) in its domain, monoconjugacy means
\[
\ker\bigl(D\operatorname{Exp}_e\vert_{tv}\bigr)\ne\{0\},
\]
and epiconjugacy means that \(D\operatorname{Exp}_e\vert_{tv}\) is not surjective. These notions can differ in infinite dimensions. Under a Jacobi-field formulation, monoconjugacy is expressed by a nonzero Jacobi field along \(\gamma_v\) vanishing at times \(0\) and \(t\), when the usual correspondence with the differential of the exponential map applies. The source does not specify which notion of conjugacy is intended. For a chosen notion, a first conjugate time means the least positive conjugate time, if such a least time exists; its existence is not assumed here.

\begin{enumerate}
\item[(a)] Determine the \emph{order of the first conjugate point} along any \(L^2\) geodesic starting at \(e\) in \(\mathscr{D}_\mu^s(M)\). The word ``order'' is retained from the source: the record does not define whether it means multiplicity, a time of occurrence, or another invariant. If multiplicity of a monoconjugate point is intended, a possible quantity is \(\dim\ker(D\operatorname{Exp}_e\vert_{tv})\), but this interpretation and its finiteness are not supplied assumptions. The question also requires a convention or restriction for geodesics having no first conjugate point.
\item[(b)] Determine whether there is a relation between the existence of conjugate points in \(\mathscr{D}_\mu^s(M)\) and Arnold's stability criterion for stationary flows on \(M\). A stationary flow means here a time-independent velocity field generating the corresponding fluid motion. The source does not specify the version or hypotheses of Arnold's criterion, the meaning of stability, the class of stationary flows, or the direction and form of the requested relation. In particular, it does not explicitly restrict the conjugate points in this part to geodesics generated by those stationary flows.
\end{enumerate}
Thus both questions are retained as source questions with unresolved qualifications, rather than strengthened into a universal existence assertion or a specified stability equivalence.
\subsection{Short English statement}
The source asks for the order of the first conjugate point along \(L^2\) geodesics in the volume-preserving diffeomorphism group, and for a connection between conjugate points and Arnold's stability criterion for stationary flows. The geometric hypotheses, the meaning of ``order'' and conjugacy, and the intended stability criterion require clarification. The supplied registry labels this problem open; that label is not independently verified.
\subsection{Sources}
\begin{itemize}
\item Local UnsolvedMath registry, permanent ID \(800817\), exact source code \texttt{catalogue entry}, title ``Geometric Hydrodynamics in Open Problems --- Problem 16.'' The supplied record attributes the problem to Khesin, Misiolek, and Modin, with proposed year \(2022\), and labels its status \texttt{open}.
\item Khesin, Misiolek, and Modin, \emph{Geometric Hydrodynamics in Open Problems} (2022), Problem 16, as identified by the supplied source record: \url{https://arxiv.org/abs/2205.01143}. The record reports extraction by \texttt{arxiv\_latex} and an access date of 2026-10-16. These are supplied provenance details; the source has not been independently visited or verified for this draft.
\end{itemize}
\end{document}
